% Created 2024-07-05 Fri 00:12
% Intended LaTeX compiler: pdflatex
\documentclass[acmsmall,nonacm]{acmart}
\usepackage[utf8]{inputenc}
\usepackage[T1]{fontenc}
%\usepackage{graphicx}
%\usepackage{grffile}
%\usepackage{longtable}
%\usepackage{wrapfig}
%\usepackage{rotating}
%\usepackage[normalem]{ulem}
\usepackage{amsmath}
\usepackage{textcomp}
\usepackage{capt-of}
\usepackage{hyperref}
\usepackage{newunicodechar}
\usepackage{pmboxdraw}
\RequirePackage{fancyvrb}
\usepackage{anyfontsize}
\usepackage{pifont}
\usepackage{fancyvrb}
\usepackage[shortcuts]{extdash}

\newenvironment{Snippet}{\Verbatim[samepage=true,fontsize=\small]}{\endVerbatim}

\DefineVerbatimEnvironment{verbatim}{Verbatim}{samepage=true}
\newunicodechar{─}{\textSFx}
\newunicodechar{│}{\textSFxi}
\newunicodechar{╭}{\textSFi}
\newunicodechar{╰}{\textSFii}
\newunicodechar{╮}{\textSFiii}
\newunicodechar{╯}{\textSFiv}
\newunicodechar{═}{\textSFxliii}
\newunicodechar{║}{\textSFxxiv}
\newunicodechar{╔}{\textSFxxxix}
\newunicodechar{╚}{\textSFxxxviii}
\newunicodechar{╗}{\textSFxxv}
\newunicodechar{╝}{\textSFxxvi}
\newunicodechar{◀}{$\triangleleft$}
\newunicodechar{▶}{$\triangleright$}
\newunicodechar{▮}{\hspace{0.6pt}|\hspace{0.6pt}}
\author{Panicz Maciej Godek}
\email{godek.maciek@gmail.com}

\date{\today}

\setcopyright{acmcopyright}
\copyrightyear{2024}
\acmYear{2024}


\acmConference[Substrates-25]{Substrates-25}{June 2025}{Prague, Czech Republic}

\setcopyright{rightsretained}

\title{The Philosophical Underpinnings of the GRASP Programming System}

\hypersetup{
 pdfauthor={Panicz Maciej Godek},
 pdftitle={The Philosophical Underpinnings of the GRASP Programming System},
 pdfkeywords={},
 pdfsubject={},
 pdfcreator={Emacs 27.1 (Org mode 9.3)}, 
 pdflang={English}}
\begin{document}
\pagenumbering{gobble}
% \tableofcontents

\begin{abstract}

\end{abstract}

\begin{CCSXML}
  <ccs2012>
  <concept>
  <concept_id>10011007.10011006.10011066.10011069</concept_id>
  <concept_desc>Software and its engineering~Integrated and visual development environments</concept_desc>
  <concept_significance>500</concept_significance>
  </concept>
  <concept>
  <concept_id>10003120.10003145.10003151.10011771</concept_id>
  <concept_desc>Human-centered computing~Visualization toolkits</concept_desc>
  <concept_significance>300</concept_significance>
  </concept>
  <concept>
  <concept>
  <concept_id>10003120.10003138.10003140</concept_id>
  <concept_desc>Human-centered computing~Ubiquitous and mobile computing systems and tools</concept_desc>
  <concept_significance>300</concept_significance>
  </concept>
  <concept_id>10010147.10010371.10010387.10010391</concept_id>
  <concept_desc>Computing methodologies~Graphics input devices</concept_desc>
  <concept_significance>100</concept_significance>
  </concept>
  </ccs2012>
\end{CCSXML}

\ccsdesc[500]{Software and its engineering~Integrated and visual development environments}
\ccsdesc[300]{Human-centered computing~Visualization toolkits}
\ccsdesc[300]{Human-centered computing~Ubiquitous and mobile computing systems and tools}
\ccsdesc[100]{Computing methodologies~Graphics input devices}

\keywords{}

\maketitle
%\onecolumn

\section{Introduction}

Every programming tool comes with its own idea of what
programming is about. The paradigm that has been dominating
the practice of computation for many decades comes with
the underlying assumption that programs are expressed
in formalized languages using text files, and that those
files are parsed and linked together according to some
specified grammar rules.

While some notable exceptions were seen since the dawn
of computation, such as Interlisp or Smalltalk, most
general-purpose programming tools that people end up
choosing are built around the notion of files encoded
in some superset of ASCII, that usually gives programmer
some limited form of expressiveness, like in the ability
to write informal comments (which includes drawing diagrams
from ASCII characters) or to indent code in some particular
ways.

Moreover, many introductory programming courses warn
their audience against using ``word processors'' such
as Microsoft Word, so that programmers aren't even allowed
to use basic typographical tools such as font color and
size for expressing their intent, even though writing
a converter that would strip all this extra information away
would be trivial.

On the other hand, programmers who might think that
they control every single byte of the source code that
they author, are often themselves puzzled by different
conventions that various operating systems use to represent
line endings in their text files (which sometimes has
real implications to the semantics of their programs).




% 


\begin{thebibliography}{12}

\bibitem{SICP}
  Harold Abelson and Gerald Jay Sussman with Julie Sussman,
  \emph{Structure and Interpretation of Computer Programs},
  Second Edition, MIT Press, 1996, ISBN 0-262-01153-0 \\
  \url{https://mitpress.mit.edu/sicp/full-text/book/book.html}
  
\bibitem{Andersen} 
 Leif Andersen, Michael Ballantyne, and Matthias
  Felleisen. 2020. \emph{Adding Interactive Visual Syntax to Textual
  Code.} Proc. ACM Program. Lang. 4, OOPSLA, Article 222 (November
  2020), 28 pages. \url{https://doi.org/10.1145/3428290}, \\
  presentation: \url{https://www.youtube.com/watch?v=8htgAxJuK5c} \\ 
  defense talk: \url{https://www.youtube.com/watch?v=l0GfMs82PvU} \\
  online IDE: \url{https://visr.pl}

\bibitem{Bochser}Michael Eisenberg, \emph{Bochser: An Integrated Scheme Programming System},
  MIT 1985, \\
  \url{https://boxer-project.github.io/boxer-literature/theses/Bochser, An Integrated Scheme Programming System (Eisenberg, MIT MSc, 1985).pdf}
  
\bibitem{Boxer}
  Andrea DiSessa, Harold Abelson,
  \emph{Boxer: A Reconstructible Computational Medium},
  MIT 1986, \url{https://web.media.mit.edu/~mres/papers/boxer.pdf}

\bibitem{Clements}
  John Clements, Matthew Flatt, and Matthias Felleisen,
  \emph{Modeling an Algebraic Stepper.} In \emph{Proc. European
  Symposium on Programming}, 2001. \\
  \url{https://www2.ccs.neu.edu/racket/pubs/esop2001-cff.pdf}
  
\bibitem{Godek2021} Panicz Maciej Godek, \emph{GRASP: A GRAphical
Scheme Programming environment (lightning talk)}, Scheme Workshop,
  2021, \url{https://www.youtube.com/watch?v=FlOghAlCDA4}

\bibitem{Godek2023} Panicz Maciej Godek. 2023. GRASP: An Extensible
  Tactile Interface for Editing S-expressions. In Proceedings of the
  16th European Lisp Symposium (ELS’23). 
  \url{https://doi.org/10.5281/zenodo.7816633} \\ recording
  available at \url{https://www.youtube.com/watch?v=TymwS5N95aY}

\bibitem{Maritime} Abraham Aaron Maritime, \emph{An Efficient
Substitution Model Scheme Evaluator}, MIT 1997 (Master thesis)\\
  \url{https://citeseerx.ist.psu.edu/document?repid=rep1&type=pdf&doi=5222083d9e0717225e3edba0cb4e85acc50c285d}

\bibitem{Meyer} Albert R. Meyer, \emph{A Substitution Model for Scheme},
  MIT 2005 (lecture notes) \\
  \url{https://courses.csail.mit.edu/6.844/spring05-6844/handouts/submodel.pdf}
  
\bibitem{Pedersen} Matt Pedersen, \emph{Implementing a simple
substitution evaluator for a Scheme-like $\lambda$-calculus language
in Scheme}, University of Nevada, Las Vegas, 2007. \\
  \url{http://www.egr.unlv.edu/~matt/teaching/CSC789/SchemeEvaluatorSubst.pdf}
  
\bibitem{jGRASP} T. Dean Hendrix, James H. Cross II, and Larry
  A. Barowski, \emph{An Extensible Framework for Providing Dynamic
  Data Structure Visualizations in a Lightweight IDE},  SIGCSE'04,
  March 3-7, 2004, Norfolk, Virginia, USA,
  \url{https://jgrasp.org/papers/sigcse2004paper.hendrix.cross.barowski.pdf}
  
\bibitem{Turbak} Franklyn Turbak, \emph{GRASP: A Visible and
Manipulable Model for Procedural Programs}, MIT 1986
  \url{https://cs.wellesley.edu/~fturbak/pubs/turbak-masters-thesis.pdf}

\end{thebibliography}


\end{document}

